
This study tests whether coefficient of variation (CV) predicts cross-benchmark ranking stability in trust evaluation. Across 10 mock trust benchmarks, CV shows weak, non-significant correlation with mean cross-benchmark Spearman $\rho$ (r=-0.486, p=0.154), failing pre-registered criteria (r<-0.5, p<0.05). This null result provides negative evidence: simple variance-based meta-features are insufficient for prospective benchmark quality assessment.

Cross-benchmark correlation analysis reveals negative correlations (e.g., FaithfulQA-FinTrust $\rho$=-0.568), suggesting trust benchmarks measure orthogonal sub-dimensions rather than a unitary ''trustworthiness'' construct. This challenges the assumption that cross-benchmark disagreement indicates measurement failure and highlights that construct validity is prerequisite for quality assessment. Before testing whether CV predicts ''cross-benchmark stability,'' we must validate what benchmarks measure via factor analysis.

This work contributes: (1) rigorous negative evidence constraining future tool development toward multi-feature models, (2) documentation of cross-benchmark correlation patterns revealing construct validity issues, and (3) a methodological framework for hypothesis-driven benchmark research with pre-registration and falsification criteria.

The critical limitation is the use of mock benchmark data. Real leaderboards have structured biases absent from synthetic data. The null result (r=-0.486) must be validated with real TrustLLM/HaluBench/TruthfulQA leaderboard data before external validity can be claimed.

A three-tier research roadmap addresses this limitation: (1) real-data replication (CRITICAL, 1 week), (2) multi-feature models combining CV, IQR, split-half reliability, ceiling proximity, and skewness (2-3 months), and (3) factor-analytic construct validation separating reliability assessment from validity assessment (6-12 months).

Simple variance metrics cannot substitute for principled construct validation. Future benchmark quality tools must be construct-aware, multi-feature, and grounded in psychometric theory, not leaderboard summary statistics in isolation. Methodology without theory generates null results; theory-driven methodology generates scientific progress, even when hypotheses are refuted.
